\subsection*{一步一步建立直观}

\paragraph{1. 先找两个倍率的中心}

二维矩阵的特征多项式是 $\lambda^2-\operatorname{tr}(A)\lambda+\det(A)$。两根之和由迹决定，所以先把未知根写成中心加减偏移。右图对称例子的中心为两个对角项的平均数，这一步减少了未知信息。



\paragraph{2. 再让乘积决定偏移大小}

两根相乘等于行列式，因此偏移的平方等于中心平方减去行列式。对实对称矩阵，这恰是 $((a-d)/2)^2+b^2$，不会为负。可用右图改变非对角项，观察中心不变而两根向两侧分开的过程。



\paragraph{3. 速度不能代替适用条件}

公式算出的根还应代回特征方程并检查方向。一般实矩阵的判别式可以为负，此时实平面中找不到非零的保持方向。快速方法的价值在于识别结构，并非免去对数域、重根和可对角化条件的判断。
